\chapter[Limits of empirical calibration]{Why empirical averages cannot certify global calibration}
\label{ch:dc-calibration-limits}

A finite record can estimate the local fluctuations it samples without
controlling fluctuations in every hidden state. The distinction matters
when an entropy certificate requires a pointwise diffusivity envelope or
a supremum bound on generator derivatives. Smoothness of each possible
model, by itself, does not provide the common numerical constants needed
for that extrapolation.

This chapter proves a precise obstruction. Over a class of smooth compact
diffusions, any uniformly covered upper confidence bound for global local
diffusivity must be infinite with probability at least its nominal coverage,
even under a benign reference process. The result holds with a bounded
detector, known uniform stationary density, fixed hidden dimension, a common
ellipticity lower bound, and a common mixing gap. It even holds when the
complete hidden trajectory is observed continuously and without noise.
Two further constructions separate the regularity assumptions: an upper
bound on diffusivity does not bound $Lf$, and upper bounds on both
diffusivity and $Lf$ do not bound $L^2f$.

These are statements about global or stationary essential suprema.
They do not deny that an ideal continuous path determines local volatility
along that path. Every model constructed below is reversible and has zero
ordinary entropy production; no change in an ideal-law entropy minimum
is asserted.

\section{A compact diffusion with an arbitrarily high rare bump}

Let $(x,z)$ lie on the torus $\mathbb T^4=(\mathbb R/2\pi\mathbb Z)^4$,
where $z$ has three coordinates. Near the origin in its fundamental cube
$[-\pi,\pi]^3$, use the Euclidean norm $|z|$. Choose a smooth radial profile
$b$ satisfying
\[
 0\leq b\leq1,\qquad b(s)=1\ (0\leq s\leq\tfrac12),\qquad
 b(s)=0\ (s\geq1).
\]
One explicit choice on $1/2<s<1$ is
\begin{equation}
 b(s)=\frac{h(1-s)}{h(1-s)+h(s-1/2)},\qquad
 h(t)=\begin{cases}e^{-1/t},&t>0,\\0,&t\leq0.\end{cases}
 \label{eq:dc-bump-profile}
\end{equation}
Together with its constant pieces this profile is $C^\infty$. Being
constant near zero also makes $b(|z|/r)$ smooth at the origin.

For $M\geq1$ and $0<r\leq1/2$, set
\begin{equation}
 \begin{split}
 d_{M,r}(z)&=1+(M-1)b(|z|/r),\\
 L_{M,r}&=d_{M,r}(z)\partial_x^2+\Delta_z,\qquad
 D_{M,r}=\operatorname{diag}(d_{M,r},1,1,1),\\
 f(x,z)&=\cos x.
 \end{split}
 \label{eq:dc-rare-model}
\end{equation}
The bump is supported strictly inside the fundamental cube, so it extends
smoothly and periodically. In It\^o form the process satisfies
\[
 dx_t=\sqrt{2d_{M,r}(z_t)}\,dW_t^x,
 \qquad dz_t=\sqrt2\,dW_t^z,
\]
with independent Brownian components and coordinates taken modulo $2\pi$.
The detector lies in $[-1,1]$.

The normalized uniform density $\pi=(2\pi)^{-4}$ is stationary. Although
$d$ depends on $z$, its only nonconstant tensor entry is the $xx$ entry,
so $L=\operatorname{div}(D\nabla)$ and $L$ is self-adjoint in $L^2(\pi)$.
Each model is smooth, connected, compact and uniformly elliptic, with the
common lower bound $D\succeq I$. Its stationary current vanishes, giving
ordinary entropy rate $\sigma=0$.

The mixing gap does not deteriorate with the bump. For every smooth
mean-zero function $g$, the torus Poincar\'e inequality gives
\begin{equation}
 -\langle g,Lg\rangle_\pi
 =\mathbb E_\pi[d|\partial_xg|^2+|\nabla_zg|^2]
 \geq\mathbb E_\pi[|\nabla g|^2]
 \geq\mathbb E_\pi[g^2].
 \label{eq:dc-common-gap}
\end{equation}
Thus the $L^2$ spectral gap is at least one, and the mean-zero semigroup
contracts at least as fast as $e^{-t}$. The construction does not rely on
metastability or a disappearing lower ellipticity constant.

For the detector in~\eqref{eq:dc-rare-model}, direct calculation yields
\begin{align}
 \Gamma_f&=\nabla f^{\mathsf T}D\nabla f=d_{M,r}\sin^2x,\\
 Lf&=-d_{M,r}\cos x,\\
 L^2f&=(d_{M,r}^2-\Delta_zd_{M,r})\cos x.
 \label{eq:dc-rare-generator}
\end{align}
On the open plateau $|z|<r/2$, the coefficient is $M$ and all its
derivatives vanish. Consequently,
\begin{equation}
 \sup\Gamma_f=M,\qquad \|Lf\|_\infty=M,\qquad
 \|L^kf\|_\infty\geq M^k\quad(k\geq1).
 \label{eq:dc-diverging-targets}
\end{equation}
These global suprema equal their stationary essential versions where
appropriate: the coefficients are continuous and $\pi$ is positive on
every open set.

The mean local diffusivity can nevertheless remain close to its reference
value. The volume of the enclosing three-dimensional ball gives
\begin{equation}
 \frac12\leq\mathbb E_\pi[\Gamma_f]
 \leq\frac12+\frac{(M-1)r^3}{12\pi^2}.
 \label{eq:dc-mean-gamma}
\end{equation}
Indeed, $\mathbb E[\sin^2x]=1/2$ and the stationary mass of the bump
support is at most $(4\pi r^3/3)/(2\pi)^3$. With $r=M^{-2}$,
the global targets diverge while the mean tends to $1/2$. This is an
explicit reason that a sample average cannot substitute for a global
envelope.

\section{An explicit bound on encountering the rare region}

\begin{lemma}[Stationary hitting bound]
\label{lem:dc-hitting}
Let $z_t$ be stationary Brownian motion with generator $\Delta$ on
$\mathbb T^3$, and let
$\tau_r=\inf\{t\geq0:|z_t|\leq r\}$. For $0<r\leq1/2$ and $T\geq0$,
\begin{equation}
 \mathbb P(\tau_r\leq T)
 \leq\min\left\{1,\ r\left(\frac9{64}+6T\right)\right\}.
 \label{eq:dc-hitting-bound}
\end{equation}
\end{lemma}

\begin{proof}
Choose a smooth radial probability density $\rho_r$ supported on
$|y|<r/2$. Its scaled Newtonian potential is
\[
 p_r(z)=r\int_{\mathbb R^3}\frac{\rho_r(y)}{|z-y|}\,dy.
\]
It is smooth and nonnegative, with
$\Delta p_r=-4\pi r\rho_r\leq0$. The radial harmonic equation and unit
total mass imply $p_r(z)=r/|z|$ for $|z|\geq r/2$.
Inside that ball the radial potential is nonincreasing with radius, so
$p_r\geq p_r(r/2)=2$. In particular, $p_r\geq1$ on $|z|\leq r$.

Take $\chi(s)=1$ for $s\leq1$, $\chi(s)=0$ for $s\geq2$, and
\[
 \chi(s)=q(s-1)\quad(1<s<2),\qquad
 q(t)=1-10t^3+15t^4-6t^5.
\]
The first two derivatives vanish at the joins, making $\chi$ a
nonnegative $C^2$ cutoff. Differentiation gives
$\max_{[0,1]}|q''|=10\sqrt3/3<6$.
Define $u_r(z)=\chi(|z|)p_r(z)$ and extend it periodically.
Its support is contained in the radius-two ball, strictly inside the
fundamental cube. Thus $u_r$ is a nonnegative $C^2$ function on the torus
and is at least one on the target ball.

For $|z|\leq1$, its Laplacian is nonpositive. On $1<|z|<2$, the radial
identity in dimension three gives
\[
 \Delta\left(\frac{\chi(s)}s\right)=\frac{\chi''(s)}s,
 \qquad s=|z|.
\]
Hence $\Delta u_r\leq6r$ everywhere; outside radius two it is zero.
The initial stationary expectation satisfies
\begin{align*}
 \mathbb E[u_r(z_0)]
 &\leq\frac r{(2\pi)^3}
   \int\rho_r(y)\int_{|z|\leq2}\frac{dz}{|z-y|}\,dy\\
 &\leq\frac r{(2\pi)^3}
      \int_{|w|\leq2+r/2}\frac{dw}{|w|}
 =\frac{r(2+r/2)^2}{4\pi^2}
 \leq\frac{81r}{64\pi^2}<\frac{9r}{64}.
\end{align*}
Here $r\leq1/2$ and $\pi>3$ supply the final constants.
Dynkin's formula at the bounded stopping time $\tau_r\wedge T$ now gives
\[
 \mathbb P(\tau_r\leq T)
 \leq\mathbb E[u_r(z_{\tau_r\wedge T})]
 \leq\mathbb E[u_r(z_0)]+6rT
 \leq r\left(\frac9{64}+6T\right).
\]
Capping by one proves~\eqref{eq:dc-hitting-bound}. The cutoff needs only
$C^2$ regularity for this argument; the diffusion coefficients in the
model remain $C^\infty$.
\end{proof}

Let $\mathbb P_0$ denote the reference model $d=1$. Couple it with
\eqref{eq:dc-rare-model} using the same uniform initial state and Brownian
drivers. Their $z$ paths agree at all times. Their $x$ coefficients agree
outside the bump, so pathwise uniqueness gives identical full hidden
trajectories until $\tau_r$. Therefore
\begin{equation}
 \operatorname{TV}\!\left(
 \mathbb P_{M,r}^{\mathrm{full},[0,T]},
 \mathbb P_0^{\mathrm{full},[0,T]}\right)
 \leq\min\left\{1,r\left(\frac9{64}+6T\right)\right\}.
 \label{eq:dc-full-path-tv}
\end{equation}
This bound is independent of the bump height $M$ and applies to complete
continuous trajectories, not merely endpoint observations.

Every common observation channel contracts total variation. Thus the same
bound applies after retaining only $f$, integrating an exposure window,
adding common independent detector noise, or applying a common randomized
processing rule. For $N$ independent stationary records with durations
summing to $T_{\rm total}$, a union coupling gives
\begin{equation}
 \operatorname{TV}(\mathbb P_{M,r}^{\rm record},\mathbb P_0^{\rm record})
 \leq\min\left\{1,r\left(\frac{9N}{64}+6T_{\rm total}\right)\right\}.
 \label{eq:dc-multiple-records}
\end{equation}
All latent intervals used by an exposure must be included in these
durations. The access contract is passive stationary observation, not an
intervention that prepares the hidden system inside a known bump.

\section{The consequence for honest upper confidence bounds}

\begin{theorem}[A global diffusivity bound cannot be learned uniformly]
\label{thm:dc-upper-confidence}
Fix a finite passive protocol covered by~\eqref{eq:dc-multiple-records}
and $0<\alpha<1$. Let $U$ be a measurable statistic of its record, including
any common external randomization, taking values in $[0,+\infty]$.
If
\begin{equation}
 \mathbb P_{M,r}\{U\geq\sup\Gamma_f\}
 =\mathbb P_{M,r}\{U\geq M\}\geq1-\alpha
 \label{eq:dc-uniform-coverage}
\end{equation}
for every $M\geq1$ and $0<r\leq1/2$, then
\begin{equation}
 \mathbb P_0\{U=+\infty\}\geq1-\alpha.
 \label{eq:dc-trivial-bound}
\end{equation}
The same conclusion holds for a uniformly covered upper bound on
$\|Lf\|_\infty$ or $\|L^2f\|_\infty$ over this class.
\end{theorem}

\begin{proof}
For each fixed $M$, total variation and
\eqref{eq:dc-uniform-coverage} give
\[
 \mathbb P_0\{U\geq M\}\geq1-\alpha
 -\min\left\{1,r\left(\frac{9N}{64}+6T_{\rm total}\right)\right\}.
\]
Let $r\downarrow0$. The resulting lower bound $1-\alpha$ holds for every
$M$. Let integer $M$ increase without bound and apply continuity from above
of probability to obtain~\eqref{eq:dc-trivial-bound}.
For $\|Lf\|_\infty$ the target equals $M$; for $\|L^2f\|_\infty$ it is
at least $M^2$. The same argument applies.
\end{proof}

In particular, a bound finite almost surely under the reference cannot
have the stated uniform coverage. The obstruction concerns the breadth
of the model class. It does not invalidate a confidence bound derived
inside a separately justified parametric family.

The theorem also applies to a claimed pointwise envelope
$\widehat a(y)$ satisfying, with uniform confidence,
\[
 \Gamma_f(x,z)\leq\widehat a(f(x,z))\quad
 \hbox{for every hidden state }(x,z).
\]
At $x=\pi/2$ and $z$ in the plateau, this condition requires
$\widehat a(0)\geq M$. Thus a uniformly covered pointwise envelope must
satisfy $\mathbb P_0\{\widehat a(0)=+\infty\}\geq1-\alpha$.
If validity is asserted only stationary-almost-everywhere, the same
conclusion holds for continuous envelopes, by continuity and positivity
of the stationary density. For arbitrary measurable almost-everywhere
envelopes, their value at the null level $y=0$ cannot be inferred this way;
an integrated formulation would be needed.

A finite stopping rule does not remove the obstruction. For a common
nonanticipating passive rule with reference stopping time $S<\infty$
almost surely, coupled stopped records agree unless $\tau_r\leq K$ or
$S>K$. Their total variation is therefore bounded by
\[
 r\left(\frac9{64}+6K\right)+\mathbb P_0(S>K),\qquad K<\infty.
\]
First let $r\downarrow0$, then $K\to\infty$. The same reference convergence
and confidence argument follow. Infinite-time observation is not included
as a finite resource.

\section{A hierarchy of missing regularity assumptions}

The preceding family allows an unbounded diffusion coefficient. The next
two variants show that supplying an upper diffusion bound still does not
supply the generator bounds used to control finite-lag approximation.

\subsection{Bounded diffusion does not bound the first generator}

Let $N$ be a positive integer multiple of four, and define
\begin{equation}
 \begin{split}
 d_{N,r}(x,z)&=1+b(|z|/r)\frac{1+\sin(Nx)}2,\\
 L_{N,r}&=\partial_x(d_{N,r}\partial_x)+\Delta_z.
 \end{split}
 \label{eq:dc-bounded-first-generator}
\end{equation}
Now $I\preceq D\preceq2I$ uniformly. The divergence-form drift
$\partial_xd_{N,r}$ is part of the generator and is essential to its
uniform stationary density. The process remains reversible, and
\eqref{eq:dc-common-gap} still proves a spectral gap at least one.
The $z$ dynamics are unchanged; outside the same rare ball, both the
$x$ drift and diffusion match the reference. Thus the full-path coupling
bounds remain valid.

For $f=\cos x$,
\[
 Lf=-d_{N,r}\cos x-\frac N2 b(|z|/r)\cos(Nx)\sin x.
\]
At $x=\pi/2$ and $z$ in the plateau, $\cos(Nx)=1$, giving $Lf=-N/2$.
Consequently,
\begin{equation}
 \Gamma_f\leq2,\qquad \|Lf\|_\infty\geq N/2.
 \label{eq:dc-bounded-first-target}
\end{equation}
Taking $N\to\infty$ and $r=N^{-2}$ makes the target diverge while the
finite-record law converges to the reference. The confidence argument
therefore remains valid for $\|Lf\|_\infty$ even when both lower and
upper diffusion bounds are known.

\subsection{Bounded diffusion and first generator do not bound the second}

Return to~\eqref{eq:dc-rare-model} and fix $M=2$. Then
$I\preceq D\preceq2I$ and
$\sup\Gamma_f=\|Lf\|_\infty=2$ for every $r$.
For the explicit profile~\eqref{eq:dc-bump-profile}, differentiating at
$s=3/4$ gives
\[
 b(3/4)=\tfrac12,\qquad b'(3/4)=-8,\qquad b''(3/4)=0.
\]
The radial Laplacian in three dimensions therefore satisfies
\[
 \Delta_3b(3/4)=b''(3/4)+\frac{2b'(3/4)}{3/4}=-\frac{64}{3}.
\]
At $x=0$ and $|z|=3r/4$, equation~\eqref{eq:dc-rare-generator} gives
\begin{equation}
 L^2f=\frac94+\frac{64}{3r^2},\qquad
 \|L^2f\|_\infty\geq\frac94+\frac{64}{3r^2}.
 \label{eq:dc-bounded-second-target}
\end{equation}
This diverges as $r\downarrow0$, while the observation law still approaches
the reference. Hence no useful uniformly finite upper confidence bound
for $\|L^2f\|_\infty$ follows merely from known bounds on $D$ and $Lf$.
Additional regularity constants are substantive assumptions, not
consequences of individual smoothness.

\section{Interpretation and reproducible checks}

The obstruction is extrapolation to hidden regions that a finite passive
record need not visit. If the bump is encountered, ideal continuous access
can reveal its increased volatility. If it is not encountered, even the
full hidden trajectory can agree exactly with the reference under the
coupling. There is no conflict with quadratic-variation reconstruction of
$\Gamma_f(X_t)$ at observed times.

The examples also clarify what can change the conclusion: a bounded
parametric family, known spatial derivative bounds, controlled preparation,
or an independently justified envelope can exclude the alternatives.
A common lower ellipticity bound, stationary density and mixing gap do
not do so. The strengthened variants retain an upper diffusion bound as
well and still separate the first and second generator requirements.
All models have $\sigma=0$; the result concerns calibration assumptions,
not hidden dissipation or a new entropy minimum.

For scale, the main family at $T=1$ and $r=M^{-2}$ has the exact analytic
record-distance upper bound $(393/64)M^{-2}$. At $M=10,100,1000$ this is
approximately $0.06140625$, $0.0006140625$ and $0.000006140625$,
respectively, while the global diffusivity supremum is $M$.
These are evaluations of the proved formulas, not simulated experiments.

The companion program implements the smooth profile and evaluates the
coupling and mean-volatility formulas. Eleven checks cover the cutoff
curvature, radial Laplacian identity, generator composition, independent
radial integration, the strengthened variants and input contracts.
Numerical outputs are non-outward-rounded diagnostics: a very small
positive mean-volatility margin can disappear under floating-point
rounding. The analytic formulas and proofs, rather than those rounded
values or a finite collection of simulations, establish the claims.
Internal independent review is distinct from external peer review or
laboratory validation; no historical priority is asserted here.
